\begin{abstract}
On-policy reward distillation converts image-level reward gradients into detached
regression targets at an intermediate state of a sampling trajectory. A design choice
that has received no attention is \emph{which} state: the strongest current method,
\opsd, builds its target at a single, low-noise query state ($\sigma\!\approx\!0.28$ on a
$10$-step schedule) and applies the loss only there. We show that this choice is tuned
for full-budget sampling and is substantially suboptimal when the sampling budget is cut.
The shortfall is not undertraining: trained to the published budget rather than a third of
it, the gap between a model's $40$-step and its $4$-step output \emph{grows}, from $18.8$
to $22.3$ in units of the measured detectable effect.

We study three answers to ``where should the target be applied''. \textbf{(A)} \emph{Move}
the query state to a higher noise level -- a one-line configuration change at zero cost.
\textbf{(B)} \emph{Shorten the training rollout}, which moves the query state implicitly.
\textbf{(C)} \emph{Propagate} the \emph{same} target to additional, higher-noise states of
the same trajectory, at no extra reward evaluation and no extra decoding.
We find that (A) and (B) are the same mechanism -- their gain curves coincide to within
measurement error -- and that both are a \emph{trade}: large gains below $10$ sampling
steps, but \emph{nothing left} at the full budget, where their change turns slightly
negative and falls below our significance threshold. Propagation (C) is not
a trade, and it is the \emph{only} one of the three that is significant in all four
prompt-set $\times$ budget cells we evaluate: it gains at every budget we measure.

We reproduce \opsd under its published protocol \emph{at its published training budget}:
all five reported evaluators land inside the agreement we measure on the untrained backbone
($1.6\%$), the largest deviation reaching $0.81$ of that tolerance. Everything else is built
on that reproduction.
On SD3.5-Medium under a matched $100$-update budget, propagation improves the reported
composite by $+5.62$ and $+3.34$ MDE at $40$ steps and by $+10.70$ and $+8.20$ MDE at
$4$ steps on two prompt sets, at $1.21\times$ training cost and with \emph{no} additional
reward evaluations. At four steps that recovers $67\%$ and $49\%$ of what tripling the
training budget buys, for a tenth of the extra compute. The noise level the supervision set reaches has an \emph{interior}
optimum: sets that stop lower or reach higher are both worse in all four
prompt-set~$\times$~budget cells, by $1.3$ to $3.2$ MDE. Propagation also composes with (B): the combination concentrates the
gain where sampling is cheapest, reaching $+16.23$ and $+15.91$ MDE at $4$ steps with
per-image win rates of $92.7\%$ and $95.4\%$ -- the largest few-step improvements we
measure. Propagation thus gives a practitioner a small family of operating points from a
single mechanism: propagate for uniform gains across budgets, and add the shorter rollout
when few-step deployment is the target.
We also map where the gain stops. It requires the published branch coefficient -- raising
$\beta$ from $0.1$ to $1.0$ reverses it, to $-15.40$ MDE at four steps -- and at the published $300$ training updates its forty-step gain is gone in both of
two seeds, while at four steps one seed keeps its lead and the other collapses after
$200$ updates. Both boundaries come from pre-registered comparisons and we report them as part of
the result. All effects
replicate across two random seeds and are calibrated against a measured
canonical-versus-canonical noise floor.
\end{abstract}
